\chapter{Koncepcja rozwiązania biometrycznego}
\label{chap:koncepcja-rozwiazania}

\section{Ogólna idea działania systemu}

Koncepcja systemu \textit{IdentityCore} opiera się na prowadzeniu lokalnego rejestru osób oraz powiązaniu każdej osoby z danymi biometrycznymi. Dane te mogą dotyczyć twarzy lub odcisku palca. Po zapisaniu osoby i jej próbek system może wykonać próbę identyfikacji na podstawie nowej próbki wejściowej.

Identyfikacja odbywa się względem danych już dostępnych w systemie. Oznacza to, że analizowana próbka nie jest oceniana samodzielnie, lecz porównywana z wcześniej zapisanymi reprezentacjami biometrycznymi osób znajdujących się w rejestrze. Wynikiem procesu jest wskazanie najlepszego dopasowania albo informacja, że dopasowanie nie zostało znalezione.

\section{Pipeline przetwarzania danych}

Pipeline systemu biometrycznego opisuje ogólny przepływ danych od momentu dostarczenia próbki wejściowej do uzyskania wyniku identyfikacji. W przypadku obu wykorzystywanych modalności -- twarzy oraz odcisku palca -- proces ma podobną strukturę koncepcyjną: dane wejściowe są przygotowywane do analizy, przekształcane do postaci reprezentacji biometrycznej, a następnie porównywane z danymi zapisanymi wcześniej w systemie.

Na podstawie porównania system wyznacza najlepiej dopasowanego kandydata i podejmuje decyzję, czy wynik można uznać za poprawne dopasowanie. Informacja o wykonanej próbie może zostać zapisana w historii, co umożliwia późniejszą analizę oraz ocenę poprawności przeprowadzonych identyfikacji.

\bigskip

Na rysunku~\ref{fig:pipeline-identitycore} przedstawiono uproszczony przepływ danych od momentu pozyskania próbki wejściowej do zapisania informacji o wykonanej próbie identyfikacji.

\begin{figure}[H]
    \centering
    \begin{tikzpicture}[
        node distance=0.6cm,
        block/.style={
            rectangle,
            draw,
            rounded corners,
            align=center,
            minimum width=9.5cm,
            minimum height=0.75cm,
            font=\small
        },
        arrow/.style={
            -{Latex[length=3mm]},
            thick
        }
    ]

        \node[block] (input) {Pozyskanie danych wejściowych\\obraz twarzy lub obraz odcisku palca};
        \node[block, below=of input] (prep) {Przygotowanie danych do analizy};
        \node[block, below=of prep] (repr) {Utworzenie reprezentacji biometrycznej};
        \node[block, below=of repr] (compare) {Porównanie z zapisanymi wzorcami};
        \node[block, below=of compare] (match) {Wyznaczenie najlepszego dopasowania};
        \node[block, below=of match] (decision) {Podjęcie decyzji o wyniku identyfikacji};
        \node[block, below=of decision] (history) {Zapis informacji o wykonanej próbie};

        \draw[arrow] (input) -- (prep);
        \draw[arrow] (prep) -- (repr);
        \draw[arrow] (repr) -- (compare);
        \draw[arrow] (compare) -- (match);
        \draw[arrow] (match) -- (decision);
        \draw[arrow] (decision) -- (history);

    \end{tikzpicture}
    \caption{Ogólny pipeline działania systemu biometrycznego IdentityCore}
    \label{fig:pipeline-identitycore}
\end{figure}

\newpage

\section{Proces rejestracji osoby}

Proces rejestracji osoby polega na utworzeniu w systemie profilu, który może zostać później wykorzystany podczas identyfikacji. Nie jest to rejestracja konta użytkownika aplikacji, lecz dodanie wpisu osoby do rejestru identyfikacyjnego wraz z przypisanymi danymi opisowymi i biometrycznymi.

W praktyce proces ten obejmuje wprowadzenie podstawowych danych osoby, dodanie zdjęcia profilowego oraz zapisanie próbek biometrycznych twarzy i odcisku palca. Dane opisowe oraz zdjęcie profilowe służą do czytelnej prezentacji osoby w aplikacji, natomiast próbki biometryczne stanowią dane odniesienia wykorzystywane podczas późniejszych prób identyfikacji.

Osoba, dla której nie zapisano danych biometrycznych danego typu, nie może zostać skutecznie zidentyfikowana z użyciem tej modalności. Oznacza to, że jakość i kompletność danych zebranych na etapie rejestracji mają bezpośredni wpływ na późniejszą skuteczność działania systemu.

\bigskip
\bigskip

Na rysunku~\ref{fig:rejestracja-osoby} przedstawiono uproszczony przebieg procesu rejestracji osoby w systemie \textit{IdentityCore}.

\begin{figure}[htbp]
    \centering
    \begin{tikzpicture}[
        node distance=0.55cm,
        block/.style={
            rectangle,
            draw,
            rounded corners,
            align=center,
            minimum width=7.2cm,
            minimum height=0.85cm,
            text width=6.8cm,
            font=\small
        },
        arrow/.style={
            -{Latex[length=2mm]},
            thick
        }
    ]

        \node[block] (start) {Wprowadzenie podstawowych danych osoby};
        \node[block, below=of start] (photo) {Dodanie zdjęcia profilowego};
        \node[block, below=of photo] (face) {Dodanie próbek twarzy};
        \node[block, below=of face] (finger) {Dodanie próbek odcisku palca};
        \node[block, below=of finger] (save) {Zapis danych w systemie};

        \draw[arrow] (start) -- (photo);
        \draw[arrow] (photo) -- (face);
        \draw[arrow] (face) -- (finger);
        \draw[arrow] (finger) -- (save);

    \end{tikzpicture}
    \caption{Proces rejestracji osoby w systemie \textit{IdentityCore}}
    \label{fig:rejestracja-osoby}
\end{figure}

\section{Proces identyfikacji osoby}

Proces identyfikacji rozpoczyna się od wyboru typu identyfikacji przez operatora. System może wykonać próbę na podstawie twarzy albo na podstawie odcisku palca. W obu przypadkach główna idea pozostaje taka sama: nowa próbka jest porównywana z danymi zapisanymi wcześniej w systemie.

W przypadku identyfikacji twarzy operator dostarcza obraz twarzy. System przygotowuje dane wejściowe do analizy, tworzy ich reprezentację biometryczną i porównuje ją z reprezentacjami zapisanymi dla osób w rejestrze. Jeżeli wynik porównania spełnia przyjęte warunki, system wskazuje najlepiej dopasowaną osobę. W przeciwnym razie informuje o braku dopasowania.

W przypadku identyfikacji odcisku palca operator dostarcza obraz odcisku. System przetwarza próbkę do postaci umożliwiającej porównanie z zapisanymi danymi biometrycznymi. Następnie wybierany jest najlepszy kandydat albo zwracana jest informacja, że nie znaleziono osoby odpowiadającej analizowanej próbce.

Po zakończeniu próby wynik może zostać zapisany w historii. Zapis obejmuje podstawowe informacje o wykonanej identyfikacji, takie jak typ użytej biometrii, wynik oraz powiązanie z rozpoznaną osobą, jeżeli dopasowanie zostało znalezione.

\newpage

\section{Schemat blokowy działania systemu}

Schemat blokowy systemu przedstawia ogólną organizację przepływu danych i decyzji w aplikacji \\\textit{IdentityCore}. W odróżnieniu od pipeline'u przetwarzania danych, schemat ten pokazuje nie tylko kolejne etapy analizy próbki, ale również powiązanie operatora, rejestru osób, danych biometrycznych, modułu identyfikacji, wyniku oraz historii prób.

Na poziomie koncepcyjnym można wyróżnić dwa główne tory działania systemu. Pierwszy dotyczy przygotowania danych odniesienia, czyli zapisania osoby w rejestrze oraz powiązania jej z danymi biometrycznymi. Drugi dotyczy wykonania próby identyfikacji, w której nowa próbka wejściowa jest porównywana z danymi zapisanymi wcześniej w systemie. Wynik próby może zostać następnie zapisany w historii.

\bigskip
\bigskip

Na rysunku~\ref{fig:schemat-blokowy-identitycore} przedstawiono uproszczony schemat blokowy działania systemu. Diagram ma charakter koncepcyjny i nie przedstawia szczegółów implementacyjnych.

\begin{figure}[htbp]
    \centering
    \begin{tikzpicture}[
        node distance=0.75cm and 0.65cm,
        block/.style={
            rectangle,
            draw,
            rounded corners,
            align=center,
            minimum width=3.45cm,
            minimum height=0.85cm,
            text width=3.25cm,
            font=\scriptsize
        },
        wideblock/.style={
            rectangle,
            draw,
            rounded corners,
            align=center,
            minimum width=4.6cm,
            minimum height=0.85cm,
            text width=4.35cm,
            font=\scriptsize
        },
        arrow/.style={
            -{Latex[length=2mm]},
            thick
        }
    ]

        % Główna oś systemu
        \node[wideblock] (operator) {Operator};
        \node[wideblock, below=of operator] (app) {Aplikacja \textit{IdentityCore}};

        % Lewa gałąź: dane odniesienia
        \node[block, below left=of app] (person) {Rejestr osób};
        \node[block, below=of person] (bio) {Dane biometryczne\\twarzy i odcisku palca};

        % Prawa gałąź: próba identyfikacji
        \node[block, below right=of app] (sample) {Nowa próbka\\wejściowa};
        \node[block, below=of sample] (module) {Moduł identyfikacji};

        % Wynik i historia
        \node[wideblock, below=3.5cm of app] (result) {Wynik identyfikacji};
        \node[wideblock, below=of result] (history) {Historia prób};

        % Połączenia główne
        \draw[arrow] (operator) -- (app);

        % Tor rejestracji danych
        \draw[arrow] (app) -- (person);
        \draw[arrow] (person) -- (bio);

        % Tor identyfikacji
        \draw[arrow] (app) -- (sample);
        \draw[arrow] (sample) -- (module);

        % Powiązanie danych odniesienia z identyfikacją
        \draw[arrow] (bio.east) -- (module.west);

        % Wynik i historia
        \draw[arrow] (module.west) -- (result.east);
        \draw[arrow] (result) -- (history);

    \end{tikzpicture}
    \caption{Schemat blokowy działania systemu \textit{IdentityCore}}
    \label{fig:schemat-blokowy-identitycore}
\end{figure}
